\documentclass[10pt]{article}
\usepackage[T1]{fontenc}
\usepackage[utf8]{inputenc}
\usepackage[french]{babel}
\usepackage{newpxtext,newpxmath}
\usepackage[a4paper,top=1.5cm,bottom=1.35cm,left=1.6cm,right=1.6cm]{geometry}
\usepackage{microtype}
\usepackage{xcolor}
\usepackage{titlesec}
\usepackage[most]{tcolorbox}
\usepackage{enumitem}
\usepackage{array}
\usepackage{booktabs}
\usepackage{fancyhdr}

\definecolor{navy}{HTML}{1F3A5F}
\definecolor{grey}{HTML}{54606E}
\definecolor{green}{HTML}{2C6E49}
\definecolor{amber}{HTML}{9C6A16}
\definecolor{red}{HTML}{A6322E}
\definecolor{plum}{HTML}{7A4EA8}
\definecolor{boxgray}{HTML}{F3F5F8}
\definecolor{boxred}{HTML}{FCF2F1}
\definecolor{boxgreen}{HTML}{F1F7F3}
\definecolor{boxplum}{HTML}{F7F3FB}
\definecolor{boxcode}{HTML}{F7F8FA}
\definecolor{boxedge}{HTML}{D9E0EA}

\setlength{\parindent}{0pt}
\setlength{\parskip}{2.6pt}
\linespread{1.0}

\pagestyle{fancy}\fancyhf{}
\lhead{\small\color{grey}Scella}
\rhead{\small\color{grey}P1 --- le profil de production --- 24 août 2026}
\cfoot{\small\color{grey}\thepage}
\renewcommand{\headrulewidth}{0.3pt}

\titleformat{\section}{\color{navy}\large\bfseries}{}{0pt}{}
\titlespacing{\section}{0pt}{11pt}{3pt}
\titleformat{\subsection}{\color{navy}\bfseries}{}{0pt}{}
\titlespacing{\subsection}{0pt}{7pt}{1.5pt}

\newtcolorbox{these}{colback=boxgray,colframe=boxedge,boxrule=0.5pt,arc=2.5pt,
  left=8pt,right=8pt,top=4.5pt,bottom=4.5pt}
\newtcolorbox{alerte}{colback=boxred,colframe=red!55,boxrule=0.8pt,arc=2.5pt,
  left=8pt,right=8pt,top=4.5pt,bottom=4.5pt}
\newtcolorbox{acquis}{colback=boxgreen,colframe=green!45,boxrule=0.8pt,arc=2.5pt,
  left=8pt,right=8pt,top=4.5pt,bottom=4.5pt}
\newtcolorbox{piste}{colback=boxplum,colframe=plum!45,boxrule=0.8pt,arc=2.5pt,
  left=8pt,right=8pt,top=4.5pt,bottom=4.5pt}
\newtcolorbox{extrait}{colback=boxcode,colframe=boxedge,boxrule=0.5pt,arc=1.5pt,
  left=7pt,right=7pt,top=3.5pt,bottom=3.5pt}

\newcommand{\best}[1]{\textcolor{green}{\textbf{#1}}}
\newcommand{\pire}[1]{\textcolor{red}{\textbf{#1}}}
\newcommand{\moyen}[1]{\textcolor{amber}{\textbf{#1}}}
\newcommand{\sz}[1]{\textcolor{grey}{\small\textsf{#1}}}
\newcommand{\code}[1]{\texttt{\small #1}}
\newcommand{\fl}{$\rightarrow$}
\newcommand{\mm}[1]{\textbf{\textcolor{navy}{#1}}}

\begin{document}

\begin{center}
{\Huge\bfseries\color{navy} Scella}\\[4pt]
{\large\color{grey}\itshape \mm{P1} --- le profil de production}\\[4pt]
{\small\color{grey} 24 août 2026 \; · \; architecture \; · \; tout ce qui suit a été
\emph{exécuté} avant d'être écrit}
\end{center}
\vspace{2pt}

\begin{these}
\textbf{\mm{P1} n'est pas « faire marcher la frontière ».} Elle marche : la configuration complète a
été montée à la main ce matin, en une douzaine de commandes --- garde franchi, \code{PgRunStore}
câblé, auth \textsc{oidc} \textsc{es}256 active, un run souverain servi en \textsc{http} 200 avec un
vrai jeton, la passerelle \textbf{tuée puis redémarrée}, et \code{verify} répondant
\code{verified:true, chain\_intact:true} sur le processus neuf. Rien n'a manqué.

\mm{P1} est donc autre chose, et de plus utile : rendre cette configuration \textbf{déclarée}
plutôt que tacite. Aujourd'hui elle n'existe que dans l'historique d'un terminal. Un profil, une
image, une charte et un contrôle qui la \emph{démarre} --- voilà le livrable.

Et l'exercice manuel a trouvé, en chemin, un défaut réel qui n'était pas au programme : le profil
\code{frontiere} livré affirme \code{frontière : SERVIE} dans une configuration où le binaire
\pire{fait exit 78}. Le test des profils ne peut pas le voir. C'est ce défaut, plus que le
packaging, qui justifie \mm{P1}.
\end{these}

\section{1 · Ce que le garde exige --- mesuré marche par marche}

\code{frontier\_boot\_guard} est une fonction pure à quatre préconditions, appelée avant tout
câblage d'état. Son contrat est écrit dans \code{main.rs} : \emph{« calculé depuis l'\textsc{intention}
env, \textsc{avant} tout câblage (\code{with\_keys}, \code{with\_coordinator}) : aucun état
forgeable / non-authentifié / volatil n'est même assemblé »}. Chaque manque \textbf{refuse de
démarrer} avec \code{exit 78} (\textsc{ex\_config}, non-retryable), jamais ne dégrade en silence.

Les quatre marches, éprouvées sur le binaire compilé \code{-{}-features postgres,oidc,pii,frontier-llm} :

\begin{center}\small
\renewcommand{\arraystretch}{1.3}
\begin{tabular}{@{}c >{\raggedright\arraybackslash}p{4.3cm} >{\raggedright\arraybackslash}p{10.2cm}@{}}
\toprule
& \textbf{Ce qui est ajouté} & \textbf{Ce que le binaire répond} \sz{(mesuré)} \\
\midrule
\textsc{e}1 & \code{FRONTIER\_ENABLED} seul & \pire{exit 78} --- « exige \code{EFE\_NANO\_AUDIT\_HMAC\_MASTER} valide ($\geq$ 32 o) --- sinon la provenance est scellée sous \code{test\_key()} \textsc{publique} $\Rightarrow$ reçus \textsc{r}4 \textsc{forgeables} » \\
\textsc{e}2 & $+$ master 32 o & \pire{exit 78} --- « exige l'auth \textsc{oidc} frontière --- sinon egress tiers \textsc{payant} sans authentification (tenant/rôles en clair) » \\
\textsc{e}3 & $+$ clé publique \textsc{es}256 & \pire{exit 78} --- « exige un store \textsc{durable} --- sinon les reçus \textsc{r}4 (la preuve) sont perdus à chaque rollout » \\
\textsc{e}4 & $+$ \code{DATABASE\_URL} & \best{préflight \textsc{ok}} --- provenance non-forgeable $+$ auth \textsc{oidc} \code{efe-llm} $+$ store durable $+$ capacité \textsc{pii} compilée \\
\bottomrule
\end{tabular}
\end{center}

\subsection{1.1 · Le détail qui fait la différence entre un garde et une décoration}

\code{key\_ok} n'est pas « la variable est présente ». C'est
\code{HmacKeySource::EnvHkdf.derive(TenantId::system()).is\_ok()} --- une \textbf{dérivation à
blanc}, c'est-à-dire exactement le calcul que fera chaque scellement. Le commentaire du code dit
pourquoi, et c'est la leçon centrale du dépôt appliquée ici : un contrôle de \emph{présence}
laisserait démarrer « vert » un master de 8 octets, puis rendrait 500 à chaque reçu. \best{Le garde
exécute le même contrat que ce qu'il protège}, plutôt que d'en vérifier un voisin plus facile.

\code{pii\_ok} suit la même honnêteté par l'autre bout : c'est
\code{cfg!(feature = "pii") || accept\_cleartext\_egress()} --- la \emph{capacité} de caviarder, ou
un aveu explicite. Le préflight refuse ensuite d'imprimer un « \textsc{ok} » lisse : sans la
capacité, il répète que \emph{« les prompts sortent \textsc{en clair} vers un tiers »}. Une
politique produit n'est pas forcée ; ce qui est exigé est que l'opérateur ait eu le \emph{moyen} de
choisir.

\code{store\_durable}, lui, n'est qu'une \textbf{intention} au moment du garde
(\code{cfg!(postgres) \&\& DATABASE\_URL.is\_some()}) --- le \code{connect}/\code{migrate} revalide
plus bas, et échoue franchement s'il ne peut pas. C'est le bon découpage : refuser tôt sur ce qu'on
peut savoir sans réseau, échouer tard et clairement sur ce qui exige une base.

\section{2 · Ce qui marche déjà --- la séquence exécutée ce matin}

\begin{acquis}
\textbf{Huit étapes, toutes vertes, sur de vrais processus.} Aucune n'est un test unitaire ; chacune
est ce qu'un opérateur ferait.
\end{acquis}

\begin{center}\small
\renewcommand{\arraystretch}{1.24}
\begin{tabular}{@{}c >{\raggedright\arraybackslash}p{7.0cm} >{\raggedright\arraybackslash}p{7.4cm}@{}}
\toprule
& \textbf{Geste} & \textbf{Résultat mesuré} \\
\midrule
1 & \code{cargo build -{}-features postgres,oidc,pii,frontier-llm} & \best{exit 0} --- les quatre features cohabitent \\
2 & un Postgres 16 réel \sz{(conteneur)} & prêt en 2 s \\
3 & une paire \textsc{ec} P-256, clé publique en \textsc{pem} & \code{openssl ecparam}, deux commandes \\
4 & démarrage avec la triade & \best{préflight \textsc{ok}}, \code{PgRunStore + PgMemoryStore câblés}, auth \textsc{oidc} active \\
5 & le rapport \mm{M0} au démarrage & provenance \best{\textsc{servie}}, store durable \best{\textsc{servie}}, authentification \best{\textsc{servie}}, frontière \moyen{\textsc{refusée}} \\
6 & \code{POST /v1/swarm/run} avec un \textsc{jwt} \textsc{es}256 & \best{\textsc{http} 200}, un \code{run\_id}, une racine \\
7 & la passerelle \textbf{tuée}, un \textbf{nouveau} processus démarré & 1 ligne en base avant, le processus remplacé après \\
8 & \code{GET /v1/swarm/verify/:id} sur le processus neuf & \best{\textsc{http} 200} --- \code{verified:true}, \code{root\_matches:true}, \code{chain\_intact:true}, \code{content\_intact:true} \\
\bottomrule
\end{tabular}
\end{center}

Deux choses méritent d'être soulignées dans cette liste, parce qu'elles ne sont pas évidentes.

\textbf{La ligne 5 est un succès, pas un demi-échec.} \code{frontière : REFUSÉE} est la bonne
réponse : aucune clé de fournisseur n'était posée, donc \code{coordinator\_from\_env()} rend
\code{None}, donc rien n'est câblé --- et \mm{M0} le dit dans son \emph{troisième état}, « fournie
mais refusée », en nommant la cause et la conséquence (\code{/v1/llm/*} répondra 501,
\code{/readyz} refusera le trafic). Un rapport qui aurait dit « \textsc{servie} » là aurait été le
mensonge que \mm{M0} existe pour rendre impossible.

\textbf{La ligne 7 est ce que le garde protège.} Le message d'\textsc{e}3 dit \emph{« les reçus
\textsc{r}4 sont perdus à chaque rollout »} ; l'étape 7 est ce rollout, joué pour de vrai. Et la
persistance Postgres est câblée \textbf{indépendamment} de la frontière --- un run souverain y va
aussi. C'est ce qui rend le critère de \mm{P1} démontrable \best{sans dépenser un centime}.

\section{3 · Le défaut trouvé en chemin --- et c'est lui qui justifie \mm{P1}}

\begin{alerte}
\textbf{Le profil \code{frontiere} livré ne peut pas démarrer, et son \code{.expected.json} affirme
qu'il sert la frontière.}

\code{profiles/frontiere.env}, tel quel, sur le binaire de production : \pire{exit 78} --- « exige
\code{EFE\_NANO\_AUDIT\_HMAC\_MASTER} valide ». Son \code{frontiere.expected.json}, lui, déclare :

\begin{center}\small
\begin{tabular}{@{}ll@{\hspace{2cm}}ll@{}}
\code{frontière} & \best{\textsc{servie}} & \code{store durable} & \pire{\textsc{dégradée}} \\
\code{provenance} & \pire{\textsc{dégradée}} & \code{authentification} & \pire{\textsc{dégradée}} \\
\end{tabular}
\end{center}

C'est un état qui \textbf{ne peut pas exister} : la frontière servie \emph{exige} les trois autres.
Le test \code{profile\_files} passe pourtant, et il a raison de passer --- il compare deux artefacts
qui ne modélisent ni l'un ni l'autre le garde.
\end{alerte}

\subsection{3.1 · Pourquoi le test ne peut pas le voir}

Le module \code{profile} porte \textbf{deux} chemins de rapport, et c'est légitime :

\begin{itemize}[leftmargin=1.2em,itemsep=1.5pt,topsep=2pt]
\item \code{Observed} --- le rapport \emph{vivant}, imprimé au démarrage. Il sait qu'aucun
      coordinator n'est câblé, et rend « \textsc{refusée} ». C'est lui qu'on a lu à l'étape 5.
\item \code{Resolved} --- le rapport \emph{dérivé d'un fichier} \code{.env}, sans processus. C'est
      celui que \code{profile\_files} compare au \code{.expected.json}. Pour la frontière, il ne lit
      que le drapeau, \code{f("EFE\_LLM\_FRONTIER\_ENABLED")}, sans jamais consulter le garde.
\end{itemize}

Le garde, lui, vit dans \code{main.rs}, sous \code{\#[cfg(feature = "frontier-llm")]}. Il n'est ni
compilé dans le build par défaut, ni atteint par le chemin \code{Resolved}. \textbf{Aucun des trois
--- profil, attendu, garde --- ne regarde les deux autres.}

\subsection{3.2 · Pourquoi cela compte ici plus qu'ailleurs}

Il faut être juste : le profil \emph{sait} qu'il est incomplet. Son en-tête énumère les secrets
« injectés au déploiement, jamais ici ». La lecture bienveillante est donc : \emph{ce}
\code{.expected.json} \emph{décrit le fichier seul, pas un démarrage}.

Mais cette lecture ne tient pas, pour une raison précise. \best{Le mot « \textsc{servie} » a un sens
unique dans tout le dépôt} --- \mm{M0} l'a défini, et \mm{M2} a ajouté « fournie mais refusée »
justement pour ne pas avoir à l'étirer. L'employer pour signifier « serait servie si les secrets
étaient là » rouvre l'écart entre le déclaré et le servi que \mm{M0} existe pour fermer. Et
l'endroit où il le rouvre est le pire possible : le seul chemin \textbf{payant} et
\textbf{non-souverain} du produit.

Un profil qui déclare la frontière et qui oublierait l'un des trois secrets ne serait rattrapé par
\emph{rien} avant le premier démarrage en production --- et il démarrerait alors en \code{exit 78},
c'est-à-dire au meilleur moment possible pour la sûreté et au pire pour la confiance.

\section{4 · L'architecture de \mm{P1} --- quatre pièces}

\subsection{4.1 · \mm{P1.a} --- le profil déclare ce qu'il ne porte PAS}

C'est la seule pièce qui invente quelque chose ; les trois autres assemblent.

Un profil \code{.env} ne peut pas porter ses trois exigences : \code{DATABASE\_URL} contient un mot
de passe, \code{EFE\_OIDC\_PUBLIC\_KEY\_PEM} est un blob \textbf{multiligne} qu'un fichier
\code{.env} ne sait pas représenter, et le master est un secret par construction. Le dépôt le sait
déjà et l'interdit : \code{profile-check.sh} \S2 échoue si un \code{.env} pose l'un des neuf noms
sensibles.

\mm{P1.a} ajoute donc la contrepartie qui manquait. Le profil déclare, dans une forme
\emph{lisible par une machine}, les variables qu'il \textbf{exige sans les porter} --- et un test
lie cette déclaration aux préconditions de \code{frontier\_boot\_guard} :

\begin{extrait}\small
\textbf{L'invariant :} pour tout profil livré dont \code{EFE\_LLM\_FRONTIER\_ENABLED=true},
l'ensemble des injections déclarées doit être \textbf{exactement} l'ensemble des préconditions du
garde --- ni moins (le profil mentirait par omission), ni plus (il exigerait ce que le binaire ne
demande pas). Le test lit les deux, et tombe si l'un bouge sans l'autre.
\end{extrait}

Ce que cela achète, précisément : le jour où une cinquième précondition sera ajoutée au garde, tout
profil frontière devient rouge \emph{immédiatement}, au lieu d'être découvert par un opérateur en
\code{exit 78}. C'est le même geste que \mm{M1} --- remplacer une vigilance par une structure --- et
c'est aussi le remède que \code{m9-check.sh} vient d'appliquer : \textbf{ne pas recopier un juge,
lier les deux à la même source}.

Et l'état déclaré du profil \code{frontiere} doit être corrigé : sans ses injections, il ne sert pas
la frontière, il la \emph{demande}. Le troisième état de \mm{M0} existe pour cela.

\subsection{4.2 · \mm{P1.b} --- l'image et la charte}

L'image du coin se construit déjà avec un argument de features : \code{docker/Dockerfile.coin}
accepte \code{CARGO\_FEATURES}, et \code{compose.yaml} s'en sert pour le profil « jonction »
(\code{reg-evidence}). \mm{P1.b} en fait une seconde cible avec
\code{postgres}, \code{oidc}, \code{pii} et \code{frontier-llm} --- \best{le build a été vérifié :
les quatre cohabitent}.

La charte \code{charts/efe-swarm-gateway} connaît déjà l'\textsc{oidc} (clé publique montée depuis
un \code{Secret} existant, jamais saisie dans les valeurs). Mesuré : elle ne connaît \pire{ni}
\code{DATABASE\_URL}, \pire{ni} \code{EFE\_LLM\_FRONTIER\_ENABLED}, \pire{ni}
\code{EFE\_PII\_DEFAULT\_POLICY} --- zéro occurrence. Trois des quatre préconditions n'ont donc
\textbf{aucun logement} dans la cible de déploiement. C'est mécanique à ajouter, et il faut le faire
avec la même discipline que l'\textsc{oidc} : la valeur vient d'un \code{Secret} référencé, jamais
d'un \code{values.yaml}.

\subsection{4.3 · \mm{P1.c} --- \code{scripts/p1-check.sh}, qui \emph{démarre} au lieu de décrire}

La règle du dépôt, apprise trois fois : \emph{une vérification qui ne lance pas ce que le client
lancerait ne vérifie pas le produit}. \code{m6-check.sh} avait lancé des tests et un exemple,
\code{m4-check.sh} avait validé un fichier compose --- les deux étaient verts pendant que le module
ne marchait pas.

\code{p1-check.sh} rejoue donc les huit étapes du \S2, plus la contre-épreuve du \S1 :

\begin{enumerate}[leftmargin=1.5em,itemsep=1.5pt,topsep=2pt]
\item \textbf{L'échelle du garde}, marche par marche : chacune des trois configurations incomplètes
      doit rendre \code{exit 78}, \emph{et le bon message}. Sans cela on ne saurait pas si le garde
      refuse pour la raison qu'il annonce.
\item Un Postgres réel, une paire \textsc{es}256 réelle, un démarrage réel.
\item Le rapport doit dire \code{provenance/store durable/authentification : SERVIE} et
      \code{frontière : REFUSÉE}. \best{Cette dernière assertion est la plus importante du script} :
      elle interdit de maquiller le résultat avec une fausse clé.
\item Un run souverain, un \code{kill}, un \textbf{nouveau} processus, un \code{verify} qui répond
      \code{verified:true}.
\item Le nettoyage, y compris en cas d'échec --- un conteneur Postgres qui survit à un script rouge
      fait passer le suivant pour de mauvaises raisons.
\end{enumerate}

\subsection{4.4 · \mm{P1.d} --- le moyen de forger un jeton}

Découvert à l'usage, et non prévu : \textbf{activer l'\textsc{oidc} change aussi le chemin
souverain}. \code{POST /v1/swarm/run} sans jeton répond \code{401 \{"error":"unauthenticated"\}}.
C'est correct --- mais un opérateur qui adopte le profil de production découvre que \emph{tout}
exige désormais un \textsc{jwt}, sans qu'aucune page ne lui dise comment en obtenir un.

\mm{P1.d} est donc une demi-journée de documentation exécutable : comment un \textsc{idp} doit
émettre le jeton (\textsc{es}256, \code{aud}, \code{tenant}, \code{roles}), et une recette locale
pour en forger un de test. Sans elle, \mm{P1} livre une passerelle que personne ne sait appeler.

\section{5 · Ce que \mm{P1} ne peut pas prouver --- et pourquoi c'est structurel}

\begin{alerte}
\textbf{Un aller-retour frontière réel est hors de portée de \mm{P1}, et pas seulement à cause du
prix.} \code{ANTHROPIC\_URL} est une \textbf{constante en dur} ---
\code{https://api.anthropic.com/v1/messages} --- employée à trois endroits, sans surcharge.
Idem pour les autres fournisseurs. Il n'existe donc \pire{aucun moyen} de pointer la frontière vers
un faux fournisseur local pour l'éprouver, ni vers un Azure OpenAI, un Bedrock, un mandataire
d'entreprise ou un modèle sur site.

Pour un produit dont la promesse est la \emph{souveraineté}, c'est une tension qu'il faut nommer :
un prospect qui exige un point de terminaison privé --- exactement le prospect que ce produit
devrait attirer --- ne peut pas être servi aujourd'hui. Ce n'est pas du ressort de \mm{P1} ; c'est
une ligne pour la feuille de route, et elle vaut probablement plus que \mm{P3}.
\end{alerte}

Deux autres limites, plus étroites mais qui mordront tôt :

\textbf{\textsc{es}256 seulement.} \code{DecodingKey::from\_ec\_pem} $+$
\code{Validation::new(Algorithm::ES256)} : une clé \textsc{rsa} est refusée. Or \textsc{rs}256 est
le défaut de Keycloak et d'Auth0. Le premier prospect qui branche son \textsc{idp} aura une chance
sur deux de tomber dessus. Et la clé est \textbf{statique} : ni \textsc{jwks}, ni découverte, donc
aucune rotation automatique.

\textbf{L'allowlist d'egress est fail-closed \emph{total}.} \code{EFE\_LLM\_TENANT\_ALLOWLIST}
absente $\Rightarrow$ registre vide $\Rightarrow$ \code{403 pour tout tenant} sur la frontière.
C'est la bonne position par défaut, mais elle signifie que même avec une vraie clé de fournisseur,
la frontière ne sert \emph{rien} tant que l'allowlist n'est pas écrite --- et ses identifiants
portent le suffixe du modèle (\code{claude\_<model>}), que le démarrage journalise justement pour
qu'on n'ait pas à le deviner.

\section{6 · Les pièges, nommés avant de les rencontrer}

\begin{center}\small
\renewcommand{\arraystretch}{1.26}
\begin{tabular}{@{}>{\raggedright\arraybackslash}p{5.2cm} >{\raggedright\arraybackslash}p{10.3cm}@{}}
\toprule
\textbf{Piège} & \textbf{Ce qu'il fait, et la parade} \\
\midrule
Le \textsc{pem} multiligne & Un \code{.env} ne le représente pas. Il doit venir d'un \code{Secret} monté ou d'une injection --- ce que la charte fait déjà bien pour l'\textsc{oidc}, et qu'il faut imiter pour \code{DATABASE\_URL} \\
La fausse clé de fournisseur & \code{ApiKey::from\_env} ne valide qu'un \emph{préfixe} (\code{sk-ant-}) sans appel réseau : une clé bidon suffirait à faire dire \code{frontière : SERVIE}. \pire{Ne jamais s'en servir} --- ce serait produire le vert malhonnête que tout le dépôt combat \\
L'\textsc{oidc} gagne le chemin souverain & 401 partout dès qu'il est activé. Attendu, mais à documenter (\mm{P1.d}) \\
Le nettoyage du contrôle & Un conteneur Postgres qui survit à un échec rend le run suivant vert pour de mauvaises raisons. \code{trap} qui \emph{compose} avec le ménage existant --- l'erreur commise dans \code{mz-check.sh}, où un second \code{trap ... EXIT} allait écraser celui qui tuait le serveur \\
\code{profile-check.sh} vieillit & Il annonce « les quatre profils » alors qu'il en existe \textbf{huit} (5 passerelle $+$ 3 console). Même famille que le compte périmé de \code{m9-check.sh} : un texte qui devient faux un module plus tard \\
\bottomrule
\end{tabular}
\end{center}

\section{7 · Le critère d'acceptation --- exécutable et non jouable}

\begin{acquis}
\code{bash scripts/p1-check.sh} sort \best{0}, et il n'y arrive qu'en ayant :

\begin{enumerate}[leftmargin=1.5em,itemsep=1pt,topsep=2pt]
\item obtenu \code{exit 78} \textbf{trois fois}, chacune avec le message de sa précondition ;
\item démarré une passerelle dont le rapport dit \code{provenance}, \code{store durable} et
      \code{authentification} \best{\textsc{servie}}, et \code{frontière} \moyen{\textsc{refusée}} ;
\item servi un run souverain en \textsc{http} 200 avec un \textsc{jwt} \textsc{es}256 ;
\item \textbf{tué} le processus et servi un \code{verify} \code{verified:true} depuis un
      \textbf{autre} processus ;
\item et laissé la machine dans l'état où il l'a trouvée.
\end{enumerate}

Plus, côté source : le test qui lie les injections déclarées d'un profil aux préconditions du
garde, et qui \textbf{tombe} si l'on en retire une.
\end{acquis}

Ce critère est \emph{non jouable} au sens du dépôt : on ne peut pas le satisfaire en écrivant un
document, ni en ajoutant une assertion complaisante. Chacun de ses cinq points exige un processus
qui démarre, un port qui répond, ou un processus qui meurt.

\section{8 · L'ordre, et ce que cela coûte}

\begin{center}\small
\renewcommand{\arraystretch}{1.26}
\begin{tabular}{@{}c >{\raggedright\arraybackslash}p{8.4cm} >{\raggedright\arraybackslash}p{4.4cm} c@{}}
\toprule
& \textbf{Pièce} & \textbf{Pourquoi dans cet ordre} & \textbf{Effort} \\
\midrule
1 & \mm{P1.c} --- le contrôle qui démarre & Il existe déjà en douze commandes ; l'écrire d'abord fige le critère avant de toucher au code & \sz{1 j} \\
2 & \mm{P1.a} --- le profil et son invariant & C'est la seule pièce qui ferme un défaut réel & \sz{1,5 j} \\
3 & \mm{P1.d} --- le jeton, documenté & Sans lui la passerelle est injoignable & \sz{0,5 j} \\
4 & \mm{P1.b} --- image et charte & Mécanique, et sans surprise mesurée & \sz{1,5 j} \\
\bottomrule
\end{tabular}
\end{center}

\begin{piste}
\textbf{Commencer par le contrôle, pas par le profil.} La séquence du \S2 a déjà tourné ; l'écrire
en script pendant qu'elle est fraîche coûte une journée et fixe le critère \emph{avant} que le code
ne bouge. C'est aussi ce qui distinguera un \mm{P1} fini d'un \mm{P1} déclaré fini : le jour où
quelqu'un ajoutera une cinquième précondition au garde, c'est ce script qui le dira.
\end{piste}

\section{Annexe · Rejouer chaque affirmation}

\renewcommand{\arraystretch}{1.22}
\begin{center}\small
\begin{tabular}{@{}>{\raggedright\arraybackslash}p{4.8cm} >{\raggedright\arraybackslash}p{11.2cm}@{}}
\toprule
\textbf{Affirmation} & \textbf{Ce qui la rejoue} \\
\midrule
Les 4 features cohabitent &
  \code{cargo build -p efe-swarm-gateway -{}-features postgres,oidc,pii,frontier-llm} \\
L'échelle du garde &
  \code{EFE\_LLM\_FRONTIER\_ENABLED=true target/debug/efe-swarm-gateway} \sz{\fl{} exit 78}, puis
  ajouter le master, puis le \textsc{pem}, puis \code{DATABASE\_URL} \\
Le profil \code{frontiere} ne démarre pas &
  \code{(set -a; . ./profiles/frontiere.env; set +a; target/debug/efe-swarm-gateway)} \sz{\fl{} exit 78} \\
\dots{} et son attendu dit le contraire &
  \code{python3 -m json.tool profiles/frontiere.expected.json} \sz{\fl{} la clé « frontière » vaut SERVIE} \\
Le rapport dit « refusée » sans clé &
  démarrer avec la triade, lire la ligne \code{frontière} du bloc « profil servi » \\
Un reçu survit au redémarrage &
  \code{POST /v1/swarm/run} \fl{} \code{kill} \fl{} redémarrer \fl{} \code{GET /v1/swarm/verify/:id} \\
L'\textsc{url} du fournisseur est en dur &
  \code{grep -n ANTHROPIC\_URL crates/efe-coordinator/src/llm/providers/anthropic.rs} \\
La charte ignore 3 préconditions sur 4 &
  \code{grep -rc 'DATABASE\_URL\textbar FRONTIER\textbar PII\_DEFAULT' charts/efe-swarm-gateway} \sz{\fl{} 0} \\
\bottomrule
\end{tabular}
\end{center}

\begin{these}
\small\textbf{Rappel de cadrage, inchangé dans toute la série.} Le reçu atteste l'\textbf{intégrité}
d'une suite de pas, jamais la \textbf{correction} d'une décision. Il est \emph{tamper-evident pour
un détenteur de clé}, pas \emph{trustless pour un tiers}. \mm{P1} ne change aucune de ces deux
phrases : il rend seulement \emph{reproductible} la configuration dans laquelle elles sont vraies.
Les \textbf{mesures} de ce document ont été exécutées avant d'être écrites et sont rejouables
ci-dessus ; les \textbf{durées} sont des estimations.
\end{these}

\end{document}
